\documentclass[12pt]{article}
\usepackage[utf8]{inputenc}
\usepackage[T1]{fontenc}
\usepackage[margin=1in]{geometry}
\usepackage{setspace}
\usepackage{xcolor}
\usepackage[colorlinks=true, urlcolor=blue!60!black, linkcolor=blue!60!black]{hyperref}
\setlength{\parindent}{2em}
\setlength{\parskip}{1.25em}
\renewcommand{\baselinestretch}{1.0}
\title{\textit{Shelley v. Kraemer}, 334 U.S. 1 (1948)}
\author{Case Brief}
\date{}
\begin{document}
\maketitle

% Case Brief: Shelley v. Kraemer, 334 U.S. 1 (1948)
% Constitutional Law Briefing Book — Chapter: Reconstruction Amendments & State Action

\section{Case Metadata}

\begin{description}
    \item[Case Name:] \textit{Shelley v. Kraemer} (consolidated with \textit{McGhee v. Sipes}, No.~87)
    \item[Citation:] \href{https://supreme.justia.com/cases/federal/us/334/1/}{334 U.S. 1 (1948)} (Docket Nos.~72 and 87; \href{https://www.oyez.org/cases/1940-1955/334us1}{Oyez Case Summary})
    \item[Court:] Supreme Court of the United States
    \item[Dates:] Argued January 15--16, 1948; Decided May 3, 1948
    \item[Chapter:] Reconstruction Amendments \& State Action
    \item[Vote:] 6--0 (unanimous among participating Justices), with 3 recusals
    \item[Opinion Author:] Chief Justice Fred M. Vinson (joined by Justices Hugo L. Black, Felix Frankfurter, William O. Douglas, Frank Murphy, and Harold H. Burton)
    \item[Recusals:] Justices Stanley F. Reed, Robert H. Jackson, and Wiley B. Rutledge took no part in the consideration or decision of the cases
    \item[Lower Court Citations:] \textit{Kraemer v. Shelley}, 355 Mo.~814, 198 S.W.2d 679 (1946) (en banc), \textit{cert. granted}, 331 U.S. 803 (1947); \textit{Sipes v. McGhee}, 316 Mich.~614, 25 N.W.2d 638 (1947), \textit{cert. granted}, 331 U.S. 804 (1947)
    \item[Disposition:] Reversed (both the judgment of the Supreme Court of Missouri and the judgment of the Supreme Court of Michigan)
    \item[Companion Case:] \href{https://supreme.justia.com/cases/federal/us/334/24/}{\textit{Hurd v. Hodge}, 334 U.S. 24 (1948)} (decided the same day, barring judicial enforcement of racially restrictive covenants in the District of Columbia under Revised Statutes \S~1978, now 42 U.S.C.~\S~1982, and federal public policy)
\end{description}

\section{Facts and Procedural History}

\textit{Shelley v. Kraemer} consolidated two state-court appeals---one from Missouri and one from Michigan---testing whether state courts could enforce private racially restrictive covenants barring non-white families from owning or occupying residential real estate (\href{https://supreme.justia.com/cases/federal/us/334/1/}{334 U.S. at 4--7}).

\subsection{The Missouri Litigation (\textit{Shelley v. Kraemer}, No.~72)}
On February 16, 1911, thirty of thirty-nine property owners fronting both sides of Labadie Avenue between Taylor Avenue and Cora Avenue in St.~Louis, Missouri, signed and subsequently recorded a fifty-year restrictive covenant covering forty-seven of the fifty-seven parcels in the designated district (\href{https://supreme.justia.com/cases/federal/us/334/1/}{334 U.S. at 4--5}). The covenant provided that, as a condition precedent to sale attaching to the land for fifty years, no part of the property could be ``occupied by any person not of the Caucasian race,'' explicitly intending to bar occupancy as owners or tenants ``by people of the Negro or Mongolian Race'' (\href{https://supreme.justia.com/cases/federal/us/334/1/}{334 U.S. at 4--5}). Notably, five parcels in the district were already owned by Black residents in 1911 (one having been occupied by Black families since 1882), and seven of the nine owners on the south side of Labadie Avenue in the immediate vicinity of the disputed parcel never signed the agreement (\href{https://supreme.justia.com/cases/federal/us/334/1/}{334 U.S. at 5}).

On August 11, 1945, petitioners J.D. and Ethel Shelley, an African-American couple, received a warranty deed for valuable consideration from Josephine Fitzgerald (an intermediary used by a real estate dealer named Bishop, who acted as the Shelleys' agent while concealing his own ownership) (\href{https://supreme.justia.com/cases/federal/us/334/1/}{334 U.S. at 5 \& n.1}). The trial court found that the Shelleys had no actual knowledge of the 1911 restrictive agreement at the time of purchase (\href{https://supreme.justia.com/cases/federal/us/334/1/}{334 U.S. at 5}).

On October 9, 1945, respondents Louis and Fern Kraemer, owners of another restricted parcel in the district, sued in the Circuit Court of the City of St.~Louis to restrain the Shelleys from taking possession and to divest title from them (\href{https://supreme.justia.com/cases/federal/us/334/1/}{334 U.S. at 6}). The trial court denied relief on state-law contract grounds, holding that the covenant was never finalized because the signatories intended it to take effect only upon obtaining the signatures of all district property owners (\href{https://supreme.justia.com/cases/federal/us/334/1/}{334 U.S. at 6}). Sitting en banc, the Supreme Court of Missouri reversed, held the covenant effective, directed the trial court to grant the requested injunctive and title-divesting relief, and ruled that judicial enforcement violated no federal constitutional right (\textit{Kraemer v. Shelley}, 355 Mo.~814, 198 S.W.2d 679 (1946); \href{https://supreme.justia.com/cases/federal/us/334/1/}{334 U.S. at 6}). By the time of the Missouri Supreme Court's decision, the Shelleys were already occupying the home.

\subsection{The Michigan Litigation (\textit{McGhee v. Sipes}, No.~87)}
In June 1934, Benjamin and Anna Ferguson executed and recorded a restrictive covenant on residential property in Detroit, Michigan, providing: ``This property shall not be used or occupied by any person or persons except those of the Caucasian race'' (\href{https://supreme.justia.com/cases/federal/us/334/1/}{334 U.S. at 6}). The restriction was set to run until January 1, 1960, conditioned on at least eighty percent participation on both sides of the street in the block---a threshold that was met (\href{https://supreme.justia.com/cases/federal/us/334/1/}{334 U.S. at 6--7}).

By deed dated November 30, 1944, petitioners Orsel and Minnie McGhee, an African-American couple, acquired title and moved into the property (\href{https://supreme.justia.com/cases/federal/us/334/1/}{334 U.S. at 7}). On January 30, 1945, neighboring property owners subject to the restriction (respondents Benjamin Sipes et al.) sued in the Circuit Court of Wayne County. Following a hearing, the trial court entered a decree ordering the McGhees to vacate the home within ninety days and permanently enjoining them from using or occupying the premises (\href{https://supreme.justia.com/cases/federal/us/334/1/}{334 U.S. at 7}). The Supreme Court of Michigan affirmed, rejecting the McGhees' Fourteenth Amendment defenses (\textit{Sipes v. McGhee}, 316 Mich.~614, 25 N.W.2d 638 (1947); \href{https://supreme.justia.com/cases/federal/us/334/1/}{334 U.S. at 7}). The U.S. Supreme Court granted certiorari in both cases and consolidated them for argument.

\section{Issue(s) Presented}

\begin{enumerate}
    \item Do private racially restrictive covenants barring persons of designated race or color from owning or occupying residential property violate the Fourteenth Amendment standing alone?
    \item Does judicial enforcement by state courts of such private racially restrictive covenants constitute ``state action'' within the meaning of Section~1 of the Fourteenth Amendment?
    \item If judicial enforcement constitutes state action, does a state court's enforcement of a racially restrictive covenant deny the excluded purchasers the equal protection of the laws guaranteed by the Fourteenth Amendment---even if state courts would likewise enforce covenants excluding white persons?
\end{enumerate}

\section{Holding}

\begin{enumerate}
    \item \textbf{Private Covenants Standing Alone:} \textbf{No.} Private agreements to exclude persons of a designated race or color from the use or occupancy of real estate, standing alone and effectuated solely through voluntary adherence, do not violate the Fourteenth Amendment (\href{https://supreme.justia.com/cases/federal/us/334/1/}{334 U.S. at 13}).
    \item \textbf{Judicial Enforcement as State Action:} \textbf{Yes.} The actions of state courts and judicial officers in their official capacities constitute action of the State within the meaning of the Fourteenth Amendment, even when they enforce private agreements pursuant to state common-law policy (\href{https://supreme.justia.com/cases/federal/us/334/1/}{334 U.S. at 14--20}).
    \item \textbf{Equal Protection Violation:} \textbf{Yes.} By granting judicial enforcement of racially restrictive covenants to divest title or enjoin occupancy by willing, financially able purchasers from willing sellers, the States of Missouri and Michigan denied petitioners the equal protection of the laws in violation of the Fourteenth Amendment (\href{https://supreme.justia.com/cases/federal/us/334/1/}{334 U.S. at 20--23}). Having resolved both cases under the Equal Protection Clause, the Court found it unnecessary to decide petitioners' Due Process and Privileges or Immunities claims (\href{https://supreme.justia.com/cases/federal/us/334/1/}{334 U.S. at 23}).
\end{enumerate}

\section{Reasoning}

Chief Justice Vinson structured the Court's unanimous (6--0) opinion in three parts, moving systematically from the substantive scope of Fourteenth Amendment property rights, to the status of state judicial action, to the constitutional invalidity of enforcing racial exclusions.

\subsection{Part I: Prior Precedent, Property Rights, and the State-Action Threshold}
First, the Court cleared away two prior restrictive-covenant decisions. In \textit{Corrigan v. Buckley}, 271 U.S. 323 (1926), the litigation originated in District of Columbia federal courts (where the Fourteenth Amendment does not apply) and challenged only the validity of the private covenant itself, not judicial enforcement under the Fifth Amendment (\href{https://supreme.justia.com/cases/federal/us/334/1/}{334 U.S. at 8--9}). In \textit{Hansberry v. Lee}, 311 U.S. 32 (1940), the Court reversed solely on procedural due process grounds (improper res judicata binding non-parties to an earlier sham class suit) without reaching the substantive constitutionality of covenant enforcement (\href{https://supreme.justia.com/cases/federal/us/334/1/}{334 U.S. at 9}).

Next, the Court emphasized that the covenants did not regulate land \textit{use} (residential occupancy remained permitted); rather, they targeted a class defined ``wholly in terms of race or color'' (\href{https://supreme.justia.com/cases/federal/us/334/1/}{334 U.S. at 10}). Equality in acquiring, owning, enjoying, and disposing of property was an essential objective of the Fourteenth Amendment's framers, reflected contemporaneously in Section~1 of the Civil Rights Act of 1866 (R.S.~\S~1978, now 42 U.S.C.~\S~1982) (\href{https://supreme.justia.com/cases/federal/us/334/1/}{334 U.S. at 10--11}). Had Missouri or Michigan imposed these occupancy bans by statute or municipal ordinance, they would be plainly unconstitutional under \href{https://supreme.justia.com/cases/federal/us/245/60/}{\textit{Buchanan v. Warley}, 245 U.S. 60 (1917)}, \textit{Harmon v. Tyler}, 273 U.S. 668 (1927), and \textit{Richmond v. Deans}, 281 U.S. 704 (1930) (\href{https://supreme.justia.com/cases/federal/us/334/1/}{334 U.S. at 11--12}).

Yet under the foundational state-action principle of the \href{https://supreme.justia.com/cases/federal/us/109/3/}{\textit{Civil Rights Cases}, 109 U.S. 3 (1883)}, Section~1 of the Fourteenth Amendment inhibits only state action, erecting ``no shield against merely private conduct, however discriminatory or wrongful'' (\href{https://supreme.justia.com/cases/federal/us/334/1/}{334 U.S. at 13}). Consequently, private racially restrictive covenants standing alone---so long as carried out through voluntary adherence---do not violate the Fourteenth Amendment.

\subsection{Part II: State Courts as State Actors}
The decisive question was whether \textit{judicial enforcement} supplied the missing state action. Reviewing three-quarters of a century of precedent, Chief Justice Vinson demonstrated that state judicial officers acting in their official capacities have always been treated as the State under the Fourteenth Amendment---citing \textit{Virginia v. Rives}, 100 U.S. 313, 318 (1880), \textit{Ex parte Virginia}, 100 U.S. 339, 347 (1880) (``A State acts by its legislative, its executive, or its judicial authorities. It can act in no other way.''), the \href{https://supreme.justia.com/cases/federal/us/109/3/}{\textit{Civil Rights Cases}, 109 U.S. at 11, 17}, \textit{Twining v. New Jersey}, 211 U.S. 78, 90--91 (1908), and \textit{Brinkerhoff-Faris Trust \& Savings Co. v. Hill}, 281 U.S. 673, 680 (1930) (\href{https://supreme.justia.com/cases/federal/us/334/1/}{334 U.S. at 14--15}).

Crucially, the Court rejected any limitation of judicial state action to procedurally unfair trials (such as jury exclusion in \textit{Strauder v. West Virginia}, 100 U.S. 303 (1880), mob-dominated trials in \textit{Moore v. Dempsey}, 261 U.S. 86 (1923), or denial of counsel in \textit{Powell v. Alabama}, 287 U.S. 45 (1932)). Judicial enforcement of a \textit{substantive state common-law rule}---even through procedurally impeccable proceedings---equally constitutes state action, as shown by First Amendment cases overturning state common-law injunctions and convictions in \textit{AFL v. Swing}, 312 U.S. 321 (1941), \textit{Cantwell v. Connecticut}, 310 U.S. 296 (1940), and \textit{Bridges v. California}, 314 U.S. 252 (1941) (\href{https://supreme.justia.com/cases/federal/us/334/1/}{334 U.S. at 16--18 \& n.19}, citing \textit{Erie R. Co. v. Tompkins}, 304 U.S. 64 (1938)).

\subsection{Part III: Active Coercive Intervention and Rejection of Respondents' Defenses}
Applying this framework, the Court held that state action existed ``in the full and complete sense of the phrase'' (\href{https://supreme.justia.com/cases/federal/us/334/1/}{334 U.S. at 19}). In both Missouri and Michigan, willing sellers had consummated sales with willing, financially able Black buyers. ``[B]ut for the active intervention of the state courts, supported by the full panoply of state power,'' the Shelleys and McGhees would have occupied their homes freely (\href{https://supreme.justia.com/cases/federal/us/334/1/}{334 U.S. at 19}). Rather than merely abstaining, the States deployed ``the full coercive power of government'' pursuant to their common-law property doctrines to deny property rights solely because of race (\href{https://supreme.justia.com/cases/federal/us/334/1/}{334 U.S. at 19--20}). That the initial discriminatory pattern originated in a private contract did not immunize the State's coercive enforcement (\href{https://supreme.justia.com/cases/federal/us/334/1/}{334 U.S. at 20}).

Finally, the Court dismantled two counterarguments advanced by respondents:
\begin{itemize}
    \item \textbf{Rejection of the ``Equal Application'' / Symmetry Argument:} Respondents argued that because state courts stood equally ready to enforce covenants excluding white persons, enforcing anti-Black covenants did not deny equal protection. The Court noted both the empirical absence of any anti-white covenant enforcement and a deeper constitutional principle: Fourteenth Amendment rights are ``personal rights'' guaranteed to the individual (\textit{McCabe v. Atchison, T. \& S.F. Ry. Co.}, 235 U.S. 151 (1914); \textit{Missouri ex rel. Gaines v. Canada}, 305 U.S. 337 (1938); \textit{Oyama v. California}, 332 U.S. 633 (1948)). ``Equal protection of the laws is not achieved through indiscriminate imposition of inequalities'' (\href{https://supreme.justia.com/cases/federal/us/334/1/}{334 U.S. at 21--22}).
    \item \textbf{Rejection of the Access-to-Courts / Property-Rights Defense:} Denying covenantors judicial enforcement did not violate \textit{their} equal protection or property rights. No individual has a constitutional right to compel state action that denies equal protection to others, and a State's power to create and enforce property interests is bounded by the Fourteenth Amendment (\href{https://supreme.justia.com/cases/federal/us/334/1/}{334 U.S. at 22}, citing \textit{Marsh v. Alabama}, 326 U.S. 501 (1946)).
\end{itemize}

\section{Concurring and Dissenting Opinions}

There were \textbf{no concurring or dissenting opinions} in \textit{Shelley v. Kraemer}. All six participating Justices joined Chief Justice Vinson's opinion in full (\href{https://supreme.justia.com/cases/federal/us/334/1/}{334 U.S. at 23}).

Two institutional features of the vote warrant note:
\begin{itemize}
    \item \textbf{The Three Recusals:} Justices Reed, Jackson, and Rutledge ``took no part in the consideration or decision of these cases'' (\href{https://supreme.justia.com/cases/federal/us/334/1/}{334 U.S. at 23}). Although the Justices issued no public explanation in the \textit{United States Reports}, historical scholarship and contemporary accounts attribute their recusals to their ownership of residential properties subject to racially restrictive covenants.
    \item \textbf{Contrast with \textit{Hurd v. Hodge}:} In the companion District of Columbia case decided the same day, \href{https://supreme.justia.com/cases/federal/us/334/24/}{\textit{Hurd v. Hodge}, 334 U.S. 24, 36 (1948)}, Justice Frankfurter wrote a brief solo concurrence resting not on R.S.~\S~1978, but on traditional equitable discretion: ``In good conscience, it cannot be `the exercise of a sound judicial discretion' by a federal court to grant the relief here asked for when the authorization of such an injunction by the States of the Union violates the Constitution.''
\end{itemize}

\section{Key Quotations}

\begin{quote}
``That Amendment erects no shield against merely private conduct, however discriminatory or wrongful.'' \\
--- \href{https://supreme.justia.com/cases/federal/us/334/1/}{\textit{Shelley v. Kraemer}, 334 U.S. 1, 13 (1948) (Justia)}
\end{quote}

\begin{quote}
``It is clear that, but for the active intervention of the state courts, supported by the full panoply of state power, petitioners would have been free to occupy the properties in question without restraint. \ldots\ Rather, these are cases in which the States have made available to such individuals the full coercive power of government to deny to petitioners, on the grounds of race or color, the enjoyment of property rights in premises which petitioners are willing and financially able to acquire and which the grantors are willing to sell.'' \\
--- \href{https://supreme.justia.com/cases/federal/us/334/1/}{\textit{Shelley v. Kraemer}, 334 U.S. 1, 19 (1948) (Justia)}
\end{quote}

\begin{quote}
``The judicial action in each case bears the clear and unmistakable imprimatur of the State. \ldots\ State action, as that phrase is understood for the purposes of the Fourteenth Amendment, refers to exertions of state power in all forms.'' \\
--- \href{https://supreme.justia.com/cases/federal/us/334/1/}{\textit{Shelley v. Kraemer}, 334 U.S. 1, 20 (1948) (Justia)}
\end{quote}

\begin{quote}
``The rights created by the first section of the Fourteenth Amendment are, by its terms, guaranteed to the individual. The rights established are personal rights. It is, therefore, no answer to these petitioners to say that the courts may also be induced to deny white persons rights of ownership and occupancy on grounds of race or color. Equal protection of the laws is not achieved through indiscriminate imposition of inequalities.'' \\
--- \href{https://supreme.justia.com/cases/federal/us/334/1/}{\textit{Shelley v. Kraemer}, 334 U.S. 1, 22 (1948) (Justia)}
\end{quote}

\begin{quote}
``The Constitution confers upon no individual the right to demand action by the State which results in the denial of equal protection of the laws to other individuals.'' \\
--- \href{https://supreme.justia.com/cases/federal/us/334/1/}{\textit{Shelley v. Kraemer}, 334 U.S. 1, 22 (1948) (Justia)}
\end{quote}

\section{Doctrinal Significance and Modern Legacy (Through 2026)}

\subsection{Immediate Progeny and Statutory Displacement in Housing}
Five years after \textit{Shelley}, the Supreme Court closed the chief remedial loophole in \href{https://supreme.justia.com/cases/federal/us/346/249/}{\textit{Barrows v. Jackson}, 346 U.S. 249 (1953)}, holding that the Fourteenth Amendment equally bars state courts from awarding money damages against a white property owner who breaches a racially restrictive covenant by selling to non-white buyers (and relaxing prudential third-party standing so the seller could assert prospective buyers' equal protection rights). Together with \href{https://supreme.justia.com/cases/federal/us/334/24/}{\textit{Hurd v. Hodge}, 334 U.S. 24 (1948)}, \textit{Shelley} and \textit{Barrows} stripped racially restrictive covenants of legal enforceability nationwide.

In 1968, federal statutory law largely superseded \textit{Shelley} as the primary tool against private housing discrimination. Congress enacted the Fair Housing Act (Title~VIII of the Civil Rights Act of 1968, 42 U.S.C.~\S~3601 \textit{et seq.}), outlawing private discrimination in housing sales, rentals, and financing. Two months later, in \href{https://supreme.justia.com/cases/federal/us/392/409/}{\textit{Jones v. Alfred H. Mayer Co.}, 392 U.S. 409 (1968)}, the Supreme Court held that 42 U.S.C.~\S~1982 (derived from the Civil Rights Act of 1866) reaches purely private racial discrimination in property transactions and is a valid exercise of Congress's Thirteenth Amendment enforcement power to abolish ``badges and incidents of slavery.'' Because \S~1982 and the Fair Housing Act reach private actors directly, modern housing plaintiffs rarely need to invoke \textit{Shelley}'s Fourteenth Amendment state-action theory.

\subsection{Doctrinal Containment: The Enduring ``\textit{Shelley} Outlier''}
Although \textit{Shelley}'s holding on racially restrictive covenants has never been questioned, the Supreme Court and lower federal courts have systematically declined to generalize its judicial-enforcement attribution rationale across constitutional law (\href{https://lawcat.berkeley.edu/record/1121131/files/fulltext.pdf}{Mark D. Rosen, \textit{Was Shelley v. Kraemer Incorrectly Decided? Some New Answers}, 95 Calif. L. Rev. 451 (2007)}):
\begin{itemize}
    \item \textbf{Limits Even in Racial Disputes:} In \href{https://supreme.justia.com/cases/federal/us/396/435/}{\textit{Evans v. Abney}, 396 U.S. 435 (1970)}, the Court distinguished \textit{Shelley} and upheld a Georgia court's application of neutral state trust-reversion law terminating a charitable park trust (returning the land to the testator's heirs) after racial segregation in the park became unconstitutional. Because closing the park imposed an equal loss on both races and reflected neutral probate rules rather than active enforcement of ongoing racial exclusion, the Court found no Equal Protection violation.
    \item \textbf{Selective Vitality in Judicial Settings:} Where state courts themselves inject or ratify racial bias inside the judicial process, \textit{Shelley}'s core insight remains vital. In \href{https://supreme.justia.com/cases/federal/us/466/429/}{\textit{Palmore v. Sidoti}, 466 U.S. 429, 431 n.1, 433 (1984)}, the Court cited \textit{Shelley} in reversing a state custody modification based on a mother's interracial remarriage, declaring that ``[p]rivate biases may be outside the reach of the law, but the law cannot, directly or indirectly, give them effect.'' Similarly, building on \href{https://supreme.justia.com/cases/federal/us/457/922/}{\textit{Lugar v. Edmondson Oil Co.}, 457 U.S. 922 (1982)} (establishing the two-part ``fair attribution'' test for private parties acting jointly with state officials), the Court held in \href{https://supreme.justia.com/cases/federal/us/500/614/}{\textit{Edmonson v. Leesville Concrete Co.}, 500 U.S. 614 (1991)}, that a private civil litigant's racially discriminatory peremptory strike constitutes state action because the courtroom and jury-selection apparatus place the government's power and prestige behind the discrimination.
    \item \textbf{Strict State-Action Boundaries Through 2026:} Over the past decade, the Supreme Court has reinforced a rigorous public/private boundary to protect private autonomy from constitutionalization. In \href{https://supreme.justia.com/cases/federal/us/587/17-1702/}{\textit{Manhattan Community Access Corp. v. Halleck}, 587 U.S. 802 (2019)}, the Court held that a private nonprofit operating public-access cable channels is not a state actor, emphasizing that the ``public function'' exception is confined to functions \textit{traditionally and exclusively} performed by government. In \href{https://supreme.justia.com/cases/federal/us/601/22-611/}{\textit{Lindke v. Freed}, 601 U.S. 187 (2024)}, the Court held unanimously that even a sitting public official's social-media blocking constitutes state action under 42 U.S.C.~\S~1983 only if the official both (1) possessed actual authority to speak on the State's behalf on the matter and (2) purported to exercise that authority in the relevant posts. Against this disciplined modern framework, \textit{Shelley} stands as a historic high-water mark of judicial state-action doctrine rather than a general rule that court enforcement of private contracts triggers constitutional scrutiny.
\end{itemize}

\subsection{Legislative Repudiation of Historic Covenants in Land Records (Through 2026)}
Because \textit{Shelley} rendered racially restrictive covenants unenforceable without physically expunging them from county land registries, thousands of deeds, plats, and subdivision declarations nationwide still contain explicit racial exclusions. To address this lingering vestige without impairing chains of title or erasing historical evidence, the Uniform Law Commission finalized the Uniform Unlawful Restrictions in Land Records Act (UURLRA) in late 2023, backed by the American Land Title Association (\href{https://www.alta.org/advocacy/advocacy-issues/discriminatory-covenants}{ALTA, \textit{Discriminatory Covenants}}). Throughout 2024--2026, a growing wave of states and the District of Columbia have enacted the UURLRA or parallel repudiation statutes (joining earlier regimes in states such as California, Maryland, Illinois, Virginia, Connecticut, Washington, and Oklahoma), enabling homeowners and homeowners' associations to record formal amendments or repudiation statements disavowing unlawful discriminatory covenants at minimal or no cost.

\section{Pedagogical Critiques and Exam Hypotheticals}

\subsection{Pedagogical Critiques}
\begin{itemize}
    \item \textbf{Collapse of the Public/Private Distinction:} If state-court enforcement of any private legal entitlement transforms the underlying private choice into state action, the state-action doctrine dissolves. Because virtually all property, contract, and probate rights rely on judicial backing, literal application of \textit{Shelley} would subject private trespass exclusions or contractual speech limits to constitutional strict scrutiny. Lower courts therefore cabin \textit{Shelley} to racial discrimination, holding---as in \textit{United Egg Producers v. Standard Brands, Inc.}, 44 F.3d 940, 943 (11th Cir. 1995)---that enforcing private speech-restricting settlements is not state action (\href{https://lawcat.berkeley.edu/record/1121131/files/fulltext.pdf}{Mark D. Rosen, 95 Calif. L. Rev. 451, 458--60 (2007)}).
    \item \textbf{Neutral Principles and Narrower Reconstructions:} Following Herbert Wechsler's critique that \textit{Shelley} lacked a generalizable neutral principle, scholars have proposed three narrower justifications: (1)~\textit{willing-party coercion} (third parties using courts to block a voluntary sale between a willing seller and willing buyer); (2)~\textit{de facto private zoning} replicating the municipal segregation invalidated in \href{https://supreme.justia.com/cases/federal/us/245/60/}{\textit{Buchanan v. Warley}, 245 U.S. 60 (1917)}; and (3)~\textit{Thirteenth Amendment statutory grounds} under 42 U.S.C.~\S~1982, the path taken in \href{https://supreme.justia.com/cases/federal/us/334/24/}{\textit{Hurd v. Hodge}, 334 U.S. 24 (1948)}, and later generalized in \href{https://supreme.justia.com/cases/federal/us/392/409/}{\textit{Jones v. Alfred H. Mayer Co.}, 392 U.S. 409 (1968)}.
    \item \textbf{State Common Law vs. Private Terms:} Chief Justice Vinson analogized \textit{Shelley} to \textit{AFL v. Swing}, 312 U.S. 321 (1941), yet critics note that in \textit{Swing} the State itself authored the substantive common-law restriction on picketing, whereas in \textit{Shelley} the discriminatory term originated solely with private covenantors.
\end{itemize}

\subsection{Exam Hypotheticals}

\subsubsection{Hypothetical 1 (``Opposite Side'' Contrast): Willing-Party Covenants vs. Personal Trespass and Neutral Trust Reversion}
\textbf{Scenario A (\textit{Shelley} Side):} A homeowners' association sues in state court to enforce a recorded neighborhood covenant barring sales to non-Christian buyers, seeking to divest title after a willing owner sells her home to a Muslim family.

\textbf{Scenario B (``Opposite Side'' Contrast):} (1)~A private homeowner bars a neighbor from a home dinner party because of the neighbor's religion and calls police to prosecute the neighbor for criminal trespass. (2)~A state probate court applies neutral reversionary trust law to close a racially restricted charitable park and return the land to the testator's heirs.

\textbf{Analysis:} Scenario~A falls squarely under \textit{Shelley} (and violates the Fair Housing Act): third parties are enlisting judicial coercion to override a voluntary transaction and enforce neighborhood-wide exclusion. Scenario~B marks the ``opposite side'' boundary: in B(1), an \textit{unwilling} owner's personal trespass exclusion reflects private associational autonomy rather than third-party coercion; in B(2), under \href{https://supreme.justia.com/cases/federal/us/396/435/}{\textit{Evans v. Abney}, 396 U.S. 435 (1970)}, neutral application of reversion rules that closes a facility to all races equally does not violate the Fourteenth Amendment.

\subsubsection{Hypothetical 2: Judicial Enforcement of Private Nondisclosure Agreements}
\textbf{Facts:} A private contractor petitions a state court to enjoin a former employee from publishing a critical op-ed in breach of a voluntary settlement's non-disparagement clause. The employee cites \textit{Shelley}, claiming judicial enforcement triggers First Amendment strict scrutiny.

\textbf{Analysis:} Although \textit{Shelley}'s broad language invites the argument, modern courts reject extending its attribution rationale to private speech contracts (\textit{United Egg Producers}), preserving private ordering under \href{https://supreme.justia.com/cases/federal/us/587/17-1702/}{\textit{Manhattan Community Access Corp. v. Halleck}, 587 U.S. 802 (2019)}, while confining judicial state action to racial discrimination or core courtroom processes such as peremptory strikes (\href{https://supreme.justia.com/cases/federal/us/500/614/}{\textit{Edmonson v. Leesville Concrete Co.}, 500 U.S. 614 (1991)}).

\end{document}
